\par\noindent\textbf{(1)} 线性动量与连续方程是 \(u_t-fv=-g\eta_x\)、\(v_t+fu=-g\eta_y\)、\(\eta_t+H(u_x+v_y)=0\)。选择经向速度恒为 \(v=0\)，其余量设为 \(e^{i(kx-\omega t)}\) 乘以 \(y\) 的包络。由东西动量与连续式分别得到 \(\omega u=kg\eta\)、\(\omega\eta=kHu\)，两者共同决定沿赤道的色散。
\par\noindent\textbf{(2)} 消去振幅得到 \(\omega^2=gHk^2\)。令 \(c=\sqrt{gH}\)，为使北半球 \(f=\beta y>0\) 与南半球同样束缚，需选择 \(\omega=ck\)、\(u=g\eta/c\) 的东向Kelvin模式。它的相速度和群速度相等 \(v_p=v_g=c\)，在常水深线性浅水近似下不展宽。
\par\noindent\textbf{(3)} 经向动量平衡 \(fu=-g\eta_y\) 与 \(f=\beta y\)、\(u=g\eta/c\) 联立得 \(\eta_y=-(\beta y/c)\eta\)。积分 \(\eta(y)=\eta(0)\exp[-\beta y^2/(2c)]\)，所以经向俘获尺度 \(L_e=\sqrt{c/\beta}\)。这不是普通边界海岸Kelvin波，而是赤道Coriolis参数变号产生的双侧束缚。
\par\noindent\textbf{(4)} 若选择反向 \(\omega=-ck\)，有 \(u=-g\eta/c\)，进而 \(\eta_y=+(\beta y/c)\eta\)，所得 \(\exp(+\beta y^2/(2c))\) 在无界经向不可归一。故不存在同型的西向赤道Kelvin束缚波。此结论适用于 \(|y|\ll R_E\) 的β平面和近似常水深，真实海洋分层、海岸与球面曲率会修改谱。
